%!TEX root = ../main.tex
% Ficha resumen exigida por la guía de estructura MCD (versión 2025-1).
% Los campos en rojo (\pendiente) deben completarse antes de la entrega.
\chapter*{Ficha resumen}
\addcontentsline{toc}{chapter}{Ficha resumen}
\begin{small}
\renewcommand{\arraystretch}{1.25}
\begin{tabularx}{\textwidth}{@{}>{\bfseries\raggedright\arraybackslash}p{4.6cm} >{\raggedright\arraybackslash}X@{}}
\toprule
Título del proyecto & Metodología basada en métodos de ciencia de datos para pronosticar SAIDI y SAIFI en redes de distribución de una empresa de energía eléctrica en el Norte de Santander \\
1. Área de trabajo & Ciencia de datos aplicada a la calidad del servicio de energía eléctrica \\
2. Tipo de proyecto & Aplicado \\
3. Estudiantes & Edwin Henao Ochoa; Cindy Lugo Rozo; Edwin Silva Salas \\
4. Correo electrónico & johanh@javerianacali.edu.co; cclugor@javerianacali.edu.co; esilvas@javerianacali.edu.co \\
5. Dirección y teléfono & \pendiente{dirección y teléfono} \\
6. Directora & Sandra Milena Ramírez Buelvas \\
7. Vinculación de la directora & Pontificia Universidad Javeriana Cali, Departamento de Ciencias Naturales y Matemáticas \\
8. Correo de la directora & smramirez@javerianacali.edu.co \\
9. Codirector & No aplica \\
10. Grupo o empresa que lo avala & \pendiente{Centrales Eléctricas de Norte de Santander S.A. E.S.P. (CENS), si la carta de autorización lo permite} \\
11. Otros grupos o empresas & No aplica \\
12. Palabras clave & SAIDI; SAIFI; calidad del servicio eléctrico; panel espacio-temporal; Random Forest; econometría espacial; CRISP-DM \\
13. Fecha de inicio & \pendiente{mes y año de inicio} \\
14. Duración estimada (meses) & \pendiente{meses} \\
15. Resumen & Véase la sección Resumen. \\
\bottomrule
\end{tabularx}
\end{small}
